% diag_lifts_loewner.tex -- round 441, DIAG_LIFTS
% (reviewer DCCCX).  Ausgang REDUZIERT.
% Builds with pdflatex.
% Companion: verify_diag_lifts.py.
% Discovery: experiments/tfpt-discovery/diag_lifts_loewner_probe.py.
% NO RH CLAIM.  Research documentation, not a theorem of RH.
\documentclass[11pt]{article}

\usepackage[a4paper,margin=2.7cm]{geometry}
\usepackage{amsmath,amssymb,amsthm}
\usepackage{microtype}
\usepackage{booktabs}
\usepackage{xcolor}
\usepackage[colorlinks=true,linkcolor=blue!50!black,urlcolor=blue!50!black]{hyperref}

\newtheorem{lemma}{Lemma}
\newtheorem{prop}{Proposition}
\theoremstyle{remark}
\newtheorem{remark}{Remark}

\newcommand{\indm}{\operatorname{ind}_{-}}

\title{\bfseries Diagonal lift of the interpolant residual\\[2mm]
\large $n_{-}(W_Y-K_{YY}^{-1})=\#\{\lambda(W^{1/2}KW^{1/2})<1\}$;
the bound $\le 1$ is not a $d_{\min}$ theorem}
\author{Stefan Hamann}
\date{August 30, 2026}

\begin{document}
\maketitle

\begin{abstract}\noindent
R439 left $D_0=W_Y-K_{YY}^{-1}$ and, at kernel dimension
$0$, the pair
$\Delta D_0\Delta=L_{-\widetilde{m}}+\mathrm{diag}(w_Y(P_Y')^2)$.
This round asks the one lemma
\texttt{DIAG\_LIFTS}: $n_{-}(D_0)\le 1$ on the selected
sequence, equivalently whether the PD diagonal lifts all
but at most one negative direction of $L_{-\widetilde{m}}$.
$L_{\widetilde{m}}$ is the dressed Cauchy Gram
$C\,\mathrm{diag}(w\pi_Y(X)^2)\,C^T$, hence always PSD;
$X$-between-$Y$ interlacing does \emph{not} change this
inertia.
The lift count is exact:
$n_{-}(W_Y-K^{-1})=\#\{\lambda(W^{1/2}KW^{1/2})<1\}$
(and at $k_{\dim}=0$,
$\#\{\lambda(D^{-1/2}LD^{-1/2})<-1\}$).
The implication $d_{\min}=\min_i w_i K_{ii}>1\Rightarrow n_{-}\le 1$
is false already on the six-node $\mathbb{Q}$-window.
The selected bound remains a census.
Ausgang \texttt{REDUZIERT}.
No RH claim.
\end{abstract}

\medskip\noindent
\textbf{Status.}\enspace
Lemmas~\ref{lem:gram}--\ref{lem:count} are proved (finite
identities).
Proposition~\ref{prop:interlace} is a named refutation.
Proposition~\ref{prop:dmin} is a named refutation.
Propositions~\ref{prop:windows}--\ref{prop:kills} are
finite machine checks.
Ausgang \texttt{REDUZIERT}.
No $L^*$ claim.  No $R^\dagger$ claim.  No RH claim.

\section{The lemma}

\begin{lemma}[DIAG\_LIFTS, two forms]
On a window let $K=K_{YY}$ be the RKHS kernel of
$P_{<d_0}$ on the $X$-measure and $W_Y=\mathrm{diag}(u^\vee|_Y)$.
The interpolant residual is $D_0=W_Y-K^{-1}$.
\texttt{DIAG\_LIFTS} asks $\indm(D_0)\le 1$ on the selected
sequence.
At $k_{\dim}=0$ this is the pair statement: the PD
diagonal $\mathrm{diag}(w_Y(P_Y')^2)$ lifts all but at most
one negative direction of $L_{-\widetilde{m}}$.
\end{lemma}

\section{Loewner signature is a Gram}

Write $\widetilde{m}(z)=\sum_a u^\vee(x_a)\,P_Y(x_a)^2/(x_a-z)$.
The Loewner matrix at $Y$ is the Cauchy Gram
\begin{equation}\label{eq:gram}
(L_{\widetilde{m}})_{ij}
=\sum_a \frac{u^\vee(x_a)\,P_Y(x_a)^2}{(x_a-y_i)(x_a-y_j)}.
\end{equation}

\begin{lemma}[Gram]\label{lem:gram}
$L_{\widetilde{m}}=C\,\mathrm{diag}(w\pi_Y(X)^2)\,C^T$
exactly over $\mathbb{Q}$ on both r431 toys.
Hence $L_{\widetilde{m}}$ is PSD of rank $\min(|X|,|Y|)$
and $L_{-\widetilde{m}}$ is NSD of the same rank.
On the five-atom toy the inertia of $L_{-\widetilde{m}}$
is $(0,2,0)$; on the six-node toy $(0,3,0)$.
\end{lemma}

\begin{prop}[Interlacing does not set $L$-inertia]\label{prop:interlace}
The five-atom toy has no $X$-node in $(\min Y,\max Y)$;
the six-node toy has one.
Both $L_{-\widetilde{m}}$ are fully ND.
Classical ``poles between nodes change Loewner inertia''
does \emph{not} apply to this source $\widetilde{m}$:
the Gram~\eqref{eq:gram} is PSD regardless of interlacing.
Fold geometry is not an $L$-inertia statement.
\end{prop}

Dropping the diagonal (the drop-$D$ mutant) restores the
full ND signature: $\indm(L_{-\widetilde{m}})=|Y|$ on both
toys.

\section{The lift count}

At $k_{\dim}=0$,
$\Delta D_0\Delta=L_{-\widetilde{m}}+D$ with
$D=\mathrm{diag}(w_Y(P_Y')^2)\succ 0$.
Sylvester gives
$\indm(D_0)=\indm(L+D)=\#\{\lambda(D^{-1/2}LD^{-1/2})<-1\}$.
On live windows $k_{\dim}\gg 0$ the pair identity fails
(relative residual $0.4$--$400$); the Cauchy--$\pi$
evaluation of $L$ overflows.
The general identity does not need the pair:

\begin{lemma}[Count]\label{lem:count}
If $W_Y\succ 0$ and $K\succ 0$ then
\begin{equation}\label{eq:count}
\indm(W_Y-K^{-1})
=\#\bigl\{\lambda\bigl(W_Y^{1/2}KW_Y^{1/2}\bigr)<1\bigr\}.
\end{equation}
At $k_{\dim}=0$ this equals the $-1$-count of the
normalised Loewner.
On both $\mathbb{Q}$-toys the two counts agree with
$\indm(D_0)$ (five-atom: $1$; six-node: $2$).
\end{lemma}

The $-1$ gap of the pair form is a $k_{\dim}=0$ object.
The live selected coordinate is the gap of the
dual-weighted CD kernel about $1$.

\section{Christoffel and $d_{\min}$}

Cauchy--Schwarz on $K$ gives $(K^{-1})_{ii}K_{ii}\ge 1$,
exactly $5592721/252000$ on the five-atom (strict).
The r408 separator is the diagonal of the weighted kernel,
$d_{\min}=\min_i w_i K_{ii}$.

\begin{prop}[$d_{\min}>1\not\Rightarrow n_{-}\le 1$]\label{prop:dmin}
On the six-node $\mathbb{Q}$-window
$d_{\min}=210157/57600>1$ and $\indm(D_0)=2$.
The implication is false as a theorem.
It remains a MAIN/adversary census:
selected and $\chi$ windows in the sealed sample have
$d_{\min}>1$; permute/jitter/scramble have $d_{\min}<1$
and many slips.
\end{prop}

$\mathrm{diag}(D_0)$ is not the carrier: MAIN $w_9$ has
no negative diagonal entry and $\indm=1$; permute has
$32$ negative diagonal entries and $\indm=20$.

\section{Windows}

\begin{prop}[Selected / kz52]\label{prop:windows}
On $w_9,k_z15,k_z18,k_z20,k_z52$ the count~\eqref{eq:count}
equals $1$.
Gaps of $W^{1/2}KW^{1/2}$ about $1$:
$w_9$ $0.857$--$1.00018$;
$k_z15$ $0.816$;
$k_z18$ $0.640$;
$k_z52$ $0.781$ ($d_{\min}=1.272$).
The r438 razor $\lambda_2(C)-1=2.2\cdot 10^{-7}$ is a
different depth; in this coordinate $k_z52$ is one slip
with a comfortable gap.
The tight selected slip in the sealed sample is $k_z22$
($0.99981$).
\end{prop}

\begin{prop}[Kills / $\chi$]\label{prop:kills}
$\Lambda$-permute: $20$ slips, $d_{\min}=0.148$.
Jitter: $19$ / $0.628$.
Scramble: $21$ / $0.367$.
No gap at $1$.
$\chi_3$-$9$: $0$ slips, $\lambda_{\min}=1.00088$ (PD,
world-separating).
$\chi_3$-$15$: one slip at $0.99816$.
Death is the side of the $1$-threshold, consistent with
r440's collar side.
\end{prop}

\section{What remains}

\begin{remark}[Remainder]
\texttt{DIAG\_LIFTS} on Selected is equivalent to:
the dual-weighted CD kernel $W_Y^{1/2}K_{YY}W_Y^{1/2}$
has at most one eigenvalue below $1$.
This is not an interlacing count of $X$ versus $Y$,
not a $d_{\min}$ theorem, and not a sign count of
$\mathrm{diag}(D_0)$.
The named remainder is off-diagonal control of $K$
(how far the weighted kernel can dip twice below $1$
while every diagonal exceeds $1$).
The six-node toy shows that this can happen off Selected.
The selected bound is a census, not a proof.
It is not a theorem of the Riemann hypothesis.
\end{remark}

Mincut unchanged (base~$4$ / refined~$5$).
No $L^*$ claim.  No $R^\dagger$ claim.  No RH claim.

\end{document}
